%! Author = ben
%! Date = 27.11.2023

% Preamble
\documentclass[../main.tex]{subfiles}

% Packages

% Document
\begin{document}
    \chapter{Ergebnissdiskussion}
    Das generelle Konzept der App SaveWorld ist sehr gut durchdacht und bietet viele Möglichkeiten.
    Jeder Benutzer kann die App auf seine Bedürfnisse anpassen und sich so ein individuelles Erlebnis schaffen.
    Dies wird vor allem durch die personalisierten Ziele und Video-Empfehlungen ermöglicht.
    Der technische Aufbau der App ist meist gut gelungen, jedoch gibt es einige Punkte, die noch verbessert werden können.
    Hierzu zählt vor allem die serverseitige Verschlüsselung der Benutzerdaten, welche in Zukunft allerdings noch hinzugefügt wird.
    \\
    Ebenfalls war geplant, dass die App über einen Hell- und Dunkelmodus verfügt, was leider aus Zeitgründen nicht mehr umgesetzt werden konnte, aber für die Zukunft geplant ist.
    \\
    Die Auswahl der Themengebiete ist sehr gut gelungen und deckt, durch die Aufteilung in Videos und Artikel, ein breites Spektrum ab.
    Die Videos sind zum Großteil sehr gut gelungen und vermitteln die Inhalte auf eine sehr verständliche Art und Weise.
    Auch die Quizze helfen dabei, das Wissen zu festigen und zu überprüfen.
    \\
    Die Artikel sind ebenfalls gut gelungen, müssen jedoch in der Zukunft noch mit Medien, wie Bildern und Videos, erweitert werden, um die Aufmerksamkeit der Benutzer zu erhöhen.
    \\
    Auch die praktische Integration der Werkzeuge, wie dem Öko-Tracker und den $CO_2$-Rechnern, ist gut gelungen und bietet den Benutzern die Möglichkeit, die Auswirkungen ihres Handelns zu verstehen.
    \\
    Das Forum fügt eine soziale Komponente hinzu, wodurch Benutzer sich gegenseitig motivieren können und gegebenenfalls Hilfe erlangen können.
    Hierbei ist die Integration der Benachrichtigungen, aber auch die Möglichkeit, Mediendateien zu versenden, sehr gut gelungen.
    \\
    Die Planung der technischen Umsetzung hat im Großen und Ganzen zwar gut funktioniert, jedoch gab es einige Probleme, bei der Auswahl der Planungssoftware.
    Zunächst haben wir mittels GitHub Issues und Milestones gearbeitet, was jedoch nicht mehr funktioniert hat, als wir ebenfalls den Inhalt der App planen wollten.
    Der Vorteil von GitHub Issues war die direkte Integration in den Quellcode, sodass Issues mit Commits verknüpft werden konnten.
    Für die Inhalte war dieses System jedoch nicht geeignet, da diese keine Integration in den Quellcode benötigen und somit die Übersichtlichkeit verloren gegangen ist.
    Letzendlich sind wir auf das Planungsprogramm Notion umgestiegen, welches uns die Möglichkeit bietet, verschiedene Ansichten zu verwenden, aber auch eine Integration von GitHub Issues ermöglicht.
    Für die Zukunft haben wir uns vorgenommen, die Planung effizienter zu gestalten, indem wir die Planung der Inhalte und der technischen Umsetzung trennen und so die Übersichtlichkeit steigern.
    \\
    Ein weiteres Problem kam bei der Auswahl der Technologien auf.
    Zunächst haben wir uns dafür entschieden, native Apps für Android und iOS zu entwickeln, was jedoch den Nachteil mitsich gebracht hätte, dass eine Web-App nur schwer umsetzbar gewesen wäre.
    Außerdem hätten wir für jede Plattform einen eigenen Code schreiben müssen, was die Entwicklung deutlich verlangsamt hätte und die Wartung erschwert hätte.
    Letzendlich haben wir uns für eine Hybrid-App mit dem Ionic Framework entschieden, wodurch wir eine Codebasis für alle Plattformen haben.
    Ebenfalls war das Design-System von Ionic nicht besonders gut in die Entwicklung integrierbar, weshalb wir uns für das Design-System Chakra-ui entschieden haben, welches deutlich personalisierbarer ist.
    \\
    Durch diesen Aufbau der App konnten wir die Entwicklung deutlich beschleunigen und den gewünschten Funktionsumfang schneller erreichen.
    \\
    Bei der Auswahl der Quellen haben wir uns vor allem auf wissenschaftliche Quellen, Bundesbehörden und seriöse Nachrichtenquellen verlassen.
    Dies war für uns sehr wichtig, da wir den Benutzern nur vertrauenswürdige Quellen zur Verfügung stellen wollen.
\end{document}
